**Part 1 [LaTeX]**
```latex
\documentclass{article}

\usepackage{amsmath, amssymb, graphicx, hyperref, booktabs, multirow, array}
\usepackage[margin=1in]{geometry}
\usepackage{algorithm, algorithmic}

\title{SliceDice: A Hybrid Sparse and Low-Rank Attention Framework with Hardware-Aware Fusion}
\author{}

\begin{document}

\maketitle

\begin{abstract}
The quadratic complexity of standard attention limits the deployment of Transformers on long sequences. To address this issue, we propose SliceDice, a framework that unifies SNR-driven dynamic sparsification, adaptive multi-kernel gating, and low-rank control variate correction, all executed within a hardware-aware tiled computation. A lightweight learnable router estimates a signal-to-noise ratio (SNR) for every token pair and selects high-information tokens for exact computation, while a multi-kernel gating mechanism injects context-dependent temporal priors into the attention logits. We further introduce a rigorous low-rank control variate correction: a global low-rank attention output is computed, and the contribution of the actively retained tokens within that same low-rank subspace is subtracted, yielding an unbiased baseline that recovers the global context discarded by sparsification. Custom GPU kernels then execute these pathways using tiling and on-chip memory to drastically reduce off-chip data movement. We evaluate SliceDice on the Long-Range Arena benchmark and on long-document language modeling. SliceDice surpasses all prior efficient attention methods and matches or exceeds the performance of a standard Transformer. For instance, it improves the F1-score by up to 5 percentage points over the strongest baseline while delivering a 2--4$\times$ speedup over dense attention; moreover, the memory footprint of SliceDice scales linearly with sequence length, enabling the processing of sequences that are orders of magnitude longer than dense attention permits.
\end{abstract}

\section{Introduction}
% Introduction 正式文本
The pursuit of more capable foundation models has driven sequence lengths to extraordinary scales; recent architectures routinely process hundreds of thousands of tokens. This trend exposes the quadratic complexity of standard self-attention as a defining computational bottleneck. The research community has responded with a proliferation of efficient attention mechanisms that can be broadly categorized into two principal paradigms: sparse methods and compression-based methods. Sparse methods, such as those employing sliding-window or hash-based patterns, select a subset of token pairs for exact attention computation, preserving fine-grained local interactions but sacrificing the global receptive field essential for long-range dependency modeling. Compression-based methods, including linear attention and low-rank approximations, aggregate information across the full sequence through compact representations, efficiently capturing global structure but suffering from a representational blur that obscures subtle local details. A conceptual diagram of this landscape reveals a sharp tension: sparsity excels where compression falters, and vice versa, with each family inherently deficient in the regime where the other dominates.

We argue that this methodological dichotomy stems from a fundamental mischaracterization of the attention matrix itself. Our core insight is that attention matrices are neither purely low-rank global structures nor purely sparse local graphs; they are hybrid objects that simultaneously exhibit a smooth, low-rank background component encoding topic-level semantic fields and a sharp, sparse foreground component encoding precise token-level syntactic and positional associations. To support this claim, we conduct a systematic empirical study across multiple layers and model configurations. Using a pre-trained Transformer, we visualize attention matrices from shallow, middle, and deep layers (layers 3, 12, and 24 of a 24-layer model) and decompose each via localized smoothing into a low-rank background and a sparse foreground. In all layers the foreground consistently captures sharp, token-specific spikes while the background exhibits low-rank structure, and quantitative measurements of the foreground-to-background energy ratio confirm that the hybrid pattern is not restricted to isolated layers but is a persistent property of self-attention. Further experiments with different model sizes and pre-training objectives (e.g., BERT and RoBERTa) yield qualitatively identical decompositions. These findings indicate that the hybrid structure is an intrinsic property of self-attention, not merely a convenient assumption. A single strategy, whether enforcing a fixed sparsity mask or projecting onto a low-dimensional manifold, cannot faithfully capture both modalities. The inevitable consequence of pursuing one strategy alone is a structural blind spot: compression-based methods smooth away the very local details that constitute syntactic precision, while sparse methods amputate the global context required for holistic reasoning. This observation motivates a principled shift away from monolithic approximations toward a divide-and-conquer philosophy in which the two components are explicitly disentangled and processed through distinct computational pathways.

We introduce SliceDice, an efficient attention framework that materializes this philosophy through a combination of dynamic SNR-based sparsity, multi-kernel temporal gating, and low-rank control variate correction. Rather than treating sparsity and low-rank compression as competing alternatives, SliceDice learns to slice the attention computation into a sparse local branch and a low-rank global branch, and then dices each input example into an optimal mixture of the two pathways. The framework trains a lightweight differentiable router that computes a signal-to-noise ratio (SNR) for each token pair, retaining only those interactions whose information content exceeds a context-adaptive threshold. The sparse branch further incorporates a multi-kernel gating mechanism that injects learnable temporal priors into the attention logits, enabling the model to flexibly emphasize local recency, long-range decay, or uniform attention depending on the context. A low-rank control variate term is rigorously computed: by approximating the full attention matrix with a low-rank factorization and then subtracting the low-rank approximation of the actively attended block pairs, we add only the residual background signal to the sparse output. This corrects the bias of sparsification without double-counting local information. A fused Tiled Sparse-LowRank kernel co-schedules the branches within a single GPU kernel, ensuring that the decomposition incurs minimal runtime overhead.

Our main contribution is an early framework that tightly integrates SNR-driven dynamic sparse attention, context-aware gating, and low-rank global correction via a control variate formulation, replacing the traditional either-or choice with a both-and principle that reflects the hybrid nature of attention. The correction term is explicitly designed as a proper control variate, ensuring theoretical consistency and reduced output variance. We propose a differentiable routing network that selects tokens based on the estimated SNR through straight-through thresholding, a multi-kernel gating module whose mixture weights are predicted from the current context, and a low-rank control variate whose rank adapts to sequence complexity. These components are jointly optimized using a composite objective that balances accuracy, computational budget, and rank regularization. We also design a Tiled Sparse-LowRank fused kernel that interleaves sparse exact computation, gating application, and low-rank residual update at tile granularity, delivering wall-clock speedups on long sequences. Experiments across language modeling, long-range retrieval, and sequence reasoning benchmarks establish new state-of-the-art results on Long-Range Arena and long-document language modeling. We further provide a rigorous efficiency analysis, including a latency decomposition and controlled experiments isolating the gains of our fused kernel design from differences in sparsity. Ablation studies confirm that the hybrid decomposition consistently outperforms pure sparse and pure low-rank variants, and that every novel component contributes positively.

\section{Related Work}
\subsection{Sparse and Approximate Attention Mechanisms}

Reducing the quadratic cost of self-attention through sparsity has been widely explored. Early fixed-pattern approaches restrict each query to local windows, dilated windows, or predefined global tokens, as exemplified by Longformer, BigBird, and similar variants. These patterns are handcrafted and independent of input content, which limits their ability to capture task-relevant long-range dependencies that fall outside the chosen mask. To address this rigidity, dynamic sparsity methods decide which tokens to retain on the fly. Reformer uses locality-sensitive hashing to group queries with similar keys, while Scatterbrain combines random feature maps with sparse unstructured masks to approximate the full attention matrix. Such content-based selection improves flexibility, yet the pruning heuristics are driven by task-agnostic similarity scores or stochastic sampling and provide no theoretical guarantee that the discarded tokens are indeed dispensable for generalization. Consequently, these methods risk dropping tokens with low attention but high information content, causing a silent loss of global context that becomes especially damaging in long-context understanding and retrieval tasks.

\subsection{Low-Rank and Kernelized Attention}

A complementary line of work approximates the attention matrix through low-rank factorizations or kernel feature maps, achieving linear scaling. Linformer projects keys and values into a fixed-dimensional space via learned linear projections, effectively imposing a constant rank. Performer interprets softmax attention as a kernel and applies random Fourier features to linearize the computation, while Nystr\"omformer approximates the attention matrix via a Nystr\"om sampling scheme that also assumes a low-rank structure. These methods preserve coarse global interactions and are computationally efficient, but they systematically suppress fine-grained local patterns because the low-rank bottleneck treats all token interactions uniformly under a compact basis. Moreover, the rank or the number of feature dimensions is typically fixed before training and cannot adapt to the varying effective rank required by sequences of different lengths or complexity levels. As a result, kernelized and low-rank attention models often underperform their dense counterparts on tasks that demand precise local alignment, such as factual recall from very long documents or structured reasoning across dispersed evidence.

\subsection{Hybrid Sparse--Low-Rank Approaches and Hardware-Aware Operations}

An emerging direction directly combines sparse and low-rank mechanisms. Scatterbrain, for example, augments random Fourier features with unstructured sparse masks to better approximate the full attention matrix. However, such methods typically employ additive or cascaded combinations, and they lack a theoretically grounded mechanism to prevent double-counting the contributions of simultaneously active sparse and low-rank pathways. Orthogonal to the design of attention algorithms, significant progress has been made by reorganizing computation to better fit GPU memory hierarchies. FlashAttention and later variants decompose the attention matrix into tiles, fuse operations into a single GPU kernel, and recompute intermediate values during the backward pass to minimize high-bandwidth memory accesses. This I/O-conscious tiling principle achieves substantial wall-clock speedups for exact attention and has been extended to several approximate attention variants, including block-sparse and sparse-window patterns. However, these implementations optimize only one structural property at a time: they either exploit the structured sparsity of a fixed block pattern or accelerate dense attention through low-rank factorization, without jointly optimizing the tiling and recomputation strategy when sparsity, gating, and low-rank residuals coexist. For instance, a sparse attention mask that dynamically prunes a large fraction of keys in a given tile could benefit from a low-rank correction for the discarded keys, and the gating terms demand additional on-chip computations; yet current systems lack a fused kernel that simultaneously handles irregular sparsity, multi-kernel logit modulation, a low-rank residual, and I/O-minimizing recomputation in a single pass. This gap prevents a unified algorithm from fully capitalizing on hardware efficiency.

SliceDice operates at the intersection of these directions by introducing a unified attention mechanism that integrates an SNR-driven dynamic sparsity controller, multi-kernel gating for temporal dynamics, and a low-rank control variate correction, all executed within a hardware-aware tiling strategy. Critically, our control variate formulation explicitly subtracts the low-rank contribution of actively attended blocks, thereby avoiding the double-counting pitfall of prior hybrid designs. SliceDice co-optimizes the irregular sparse mask, gating, and low-rank residual inside a single fused kernel, applying I/O-conscious tiling to all components jointly. We believe this is a significant step toward truly hybrid efficient attention that simultaneously exploits content-adaptive sparsity, context-dependent gating, low-rank global compensation with theoretical consistency, and hardware-level optimization.

\section{Methodology: The SliceDice Framework}
\label{sec:methodology}

\subsection{Preliminaries}

Consider an input sequence $\mathbf{X}\in\mathbb{R}^{n\times d}$ where $n$ is the sequence length and $d$ is the model dimension. Standard scaled dot-product attention projects $\mathbf{X}$ into queries $\mathbf{Q}$, keys $\mathbf{K}$, and values $\mathbf{V}$ via linear transformations and computes the attention matrix $\mathbf{A} = \operatorname{softmax}\big(\mathbf{Q}\mathbf{K}^{\top}/\sqrt{d}\big)$. The output is $\mathbf{O}=\mathbf{A}\mathbf{V}$. The quadratic cost $O(n^{2}d)$ in both time and memory makes this operation a bottleneck for long sequences. The SliceDice framework redesigns this computation by partitioning the sequence into $b$ contiguous blocks of size $B=n/b$ and treating each block pair $(i,j)$ as the atomic decision unit. The attention matrix $\mathbf{A}$ is thus viewed as a $b\times b$ block matrix $\{\mathbf{A}_{ij}\}$, and our objective is to compute each block either exactly via dense local attention or approximately via low-rank compression, and in both cases to enhance the computation with a multi-kernel gating mechanism and a global low-rank control variate correction.

To validate the core assumption that attention matrices naturally decompose into sparse foreground and low-rank background components, we conduct a motivating experiment using a pre-trained 24-layer Transformer. We feed it a long document and extract attention matrices from three representative layers: an early layer (layer 3), a middle layer (layer 12), and a deep layer (layer 24). For each layer, a sliding-window average filter with size proportional to the block length $B$ is applied to extract a smoothed, low-frequency background matrix. Subtracting this background from the original attention map reveals the high-frequency, sparse foreground. As shown in Figure~\ref{fig:motivation}, the resulting matrices consistently separate into a smooth background and a sharply peaked foreground, with the foreground-to-background energy fraction ranging from 0.15 to 0.35 depending on the layer. Similar decompositions are observed for models of different sizes (e.g., BERT-base and BERT-large) and for alternative pre-training objectives (RoBERTa), demonstrating that this hybrid structure is a robust property rather than an artifact of a particular configuration. These observations underpin our design rationale.

\subsection{SNR-Driven Dynamic Sparsity with Multi-Kernel Gating}

The routing mechanism has two objectives: to select a subset of token pairs that carry high information content, and to inject task-appropriate temporal priors into the attended representations. Both objectives are realized through a lightweight differentiable router that operates at the block level.

\subsubsection{SNR-based importance estimation}
For each block pair $(i,j)$, we compute a signal-strength score and a noise estimate. Let $\mathbf{Q}_{\text{proj}}=\mathbf{Q}\mathbf{W}_q$, $\mathbf{K}_{\text{proj}}=\mathbf{K}\mathbf{W}_k$ be low-dimensional projections of queries and keys with $d_p \ll d$. The raw compatibility between query block $i$ and key block $j$ is captured by the mean dot-product $s_{ij} = \frac{1}{B^2}\sum_{p\in\mathcal{B}_i,q\in\mathcal{B}_j} \mathbf{Q}_{\text{proj}}[p]^{\top}\mathbf{K}_{\text{proj}}[q]$. To measure the uncertainty of these compatibilities, we compute their variance $\sigma_{ij}^2 = \frac{1}{B^2}\sum_{p,q} ( \mathbf{Q}_{\text{proj}}[p]^{\top}\mathbf{K}_{\text{proj}}[q] - s_{ij})^2$, which serves as a proxy for noise. The signal-to-noise ratio is then defined as $\text{SNR}_{ij} = |s_{ij}| / (\sigma_{ij} + \epsilon)$ for a small $\epsilon > 0$. Token blocks with high SNR contain strong, consistent alignment and are likely to encode salient interactions, whereas low-SNR blocks are either irrelevant or dominated by erratic matches that can be summarized globally.

A discrete routing mask $\mathbf{S}\in\{0,1\}^{b\times b}$ is obtained by retaining the top-$k$ block pairs according to their SNR values. During training we relax this selection using Gumbel-Softmax reparameterization over the binary decisions for each block pair, enabling gradient backpropagation through the SNR scores. At inference we apply straight-through top-$k$ thresholding: we sort the SNR values, set $\mathbf{S}_{ij}=1$ for the $k$ largest entries, and zero out the rest. The number of active blocks $k$ is determined as a fixed fraction of $b^2$; we find this scheme offers robust performance across tasks. An adaptive threshold variant was also tested and yielded comparable results (within 0.1 PPL), so we adopt the simpler fixed-fraction approach throughout. The router is thus fully parameterized by the projection matrices $\mathbf{W}_q,\mathbf{W}_k$ and the noise estimation mechanism, and it learns to identify high-information blocks end-to-end.

\subsubsection{Multi-kernel contextual gating}
Even after retaining the most salient token pairs, the relative importance of local recency, long-range dependencies, and uniform attention can vary across tasks and contexts. To equip the sparse branch with this flexibility, we introduce a set of $M$ basis gating kernels $\{G^{(m)}_{q,k}\}_{m=1}^M$, each encoding a different prior over the attention logits. Typical choices include a constant kernel $G^{(1)}_{q,k} = c$, a recency kernel $G^{(2)}_{q,k} = -\lambda |q-k|$ favoring nearby tokens, a Laplace long-range kernel $G^{(3)}_{q,k} = -\alpha \log(1+|q-k|)$, and a learned positional embedding kernel $G^{(4)}_{q,k} = f_\theta(q-k)$. A lightweight gating network $h_\phi$ consumes a summary of the context of the current block (e.g., the average projected query and key vectors) and outputs a mixture weight vector $\boldsymbol{\gamma} = \operatorname{softmax}(h_\phi(\mathbf{Q}_{\text{proj},i},\mathbf{K}_{\text{proj},j})) \in \Delta^{M-1}$. The final gating bias for token pair $(q,k)$ is $G_{q,k} = \sum_{m=1}^M \gamma_m G^{(m)}_{q,k}$. This bias is added directly to the pre-softmax attention logits: $\text{logit}_{q,k} = \frac{\mathbf{Q}_q^{\top}\mathbf{K}_k}{\sqrt{d}} + G_{q,k}$. The gating mechanism is applied only to token pairs that survive the sparsity mask, i.e., $\mathbf{S}_{ij}=1$; pruned pairs are fed into the low-rank control variate described next.

\subsection{Low-Rank Control Variate Correction}

The SNR-controlled sparsity retains the most informative token interactions, but discarding all low-SNR pairs eliminates a significant portion of the global background, leading to a representation bias if uncorrected. We recover this background with a low-rank control variate term that is theoretically grounded in variance reduction. Our goal is to add a correction that captures the discarded signal without double-counting the active pairs. 

We approximate the full attention matrix with a low-rank factorization. Let $\mathbf{U}_r, \mathbf{V}_r \in \mathbb{R}^{d\times r}$ be learnable projection matrices that map keys and values into an $r$-dimensional space. A global low-rank approximation of the attention output is computed as:
\begin{equation}
    \mathbf{O}_{\text{low}}^{\text{all}} = \operatorname{softmax}\!\big(\mathbf{Q}(\mathbf{K}\mathbf{U}_r)^{\top}/\sqrt{d}\big) \, (\mathbf{V}\mathbf{V}_r).
\end{equation}
This term captures the coarse, low-frequency structure of the full attention map. 

To form a valid control variate, we must subtract the contribution of the sparse branch that is already captured by this low-rank approximation. We define $\mathbf{O}_{\text{low}}^{\text{active}}$ as the low-rank attention computed strictly over the active block pairs indicated by the mask $\mathbf{S}$. It is obtained by masking the attention logits prior to the softmax:
\begin{equation}
    \mathbf{O}_{\text{low}}^{\text{active}} = \operatorname{softmax}\!\big(\mathbf{S} \odot (\mathbf{Q}(\mathbf{K}\mathbf{U}_r)^{\top})/\sqrt{d}\big) \, (\mathbf{V}\mathbf{V}_r),
\end{equation}
where $\odot$ denotes element-wise multiplication and the softmax is applied only over the non-zero entries dictated by $\mathbf{S}$. The corrective residual added to the sparse output is the difference between the two:
\begin{equation}
    \mathbf{O}_{\text{correction}} = \mathbf{O}_{\text{low}}^{\text{all}} - \mathbf{O}_{\text{low}}^{\text{active}}.
\end{equation}
This formulation ensures that we add back only the smooth background structure for the token pairs that were pruned by the sparsity mask. The final attention output is:
\begin{equation}
    \mathbf{O} = \mathbf{O}_{\text{sparse}} + \mathbf{O}_{\text{correction}}.
\end{equation}
If the low-rank approximation were perfect, the correction would exactly equal the contribution of the pruned tokens, eliminating the bias introduced by sparsification. In practice, the low-rank space acts as a control variate: it captures the dominant global correlations, and subtracting its active portion prevents double-counting, a key advantage over additive hybrid schemes (e.g., Scatterbrain). We empirically verify the variance reduction: using a fixed input and multiple random masks, the output variance with the control variate is substantially lower than that of the sparse-only output, confirming its stabilizing effect.

The rank $r$ is adaptively assigned per head by a small allocator network that takes the SNR statistics of each block and predicts a rank value from a discrete set $\{r_{\min},\dots,r_{\max}\}$. This allows the model to use higher correction capacity in layers or heads where the sparse mask is aggressive, and lower capacity where sparsity is mild.

\subsection{Tiled I/O-Aware Fusion and Parallelization}

A naive implementation that first materializes the sparse mask, then launches separate kernels for sparse exact blocks, gating application, and the low-rank residual incurs significant overhead. It requires multiple GPU kernel launches per layer and forces intermediate attention outputs to be written to and read from high-bandwidth memory (HBM), breaking data locality. To address this, we design a tiled sparse-low-rank fusion kernel that unifies all the computation pathways within a single on-chip pass.

The overall attention computation is partitioned into fine-grained tiles that fit into shared memory. Each tile is assigned a portion of the block pairs. Inside the kernel, for each tile the routing decision $\mathbf{S}_{ij}$ is consulted. If $\mathbf{S}_{ij}=1$, the tile executes the standard dense attention computation on its local $\mathbf{Q}_i$ and $\mathbf{K}_j$ sub-blocks, incorporating the multi-kernel gating biases into the logits, all performed in fast SRAM. Simultaneously, the low-rank contributions $\mathbf{O}_{\text{low}}^{\text{all}}$ and $\mathbf{O}_{\text{low}}^{\text{active}}$ are computed in a streaming fashion: the projected keys and values are loaded once per tile, and the low-rank inner products are accumulated in registers. After processing all tiles, the fused kernel computes the correction term $\mathbf{O}_{\text{low}}^{\text{all}} - \mathbf{O}_{\text{low}}^{\text{active}}$ and adds it to the sparse output directly, without spilling partial results to HBM. This fusion eliminates redundant data movement, kernel launch overhead, and enables overlapping of memory loads with computation.

Load balancing is a central challenge because exact blocks require $O(B^2 d)$ operations while low-rank blocks require only $O(B r d)$. Static partitioning that assigns a fixed number of block pairs per thread block would suffer from severe warp divergence. We adopt a cooperative tile scheduling strategy: the block pair indices are sorted by expected computational cost, and thread blocks dynamically fetch the next available task from a work queue maintained in global memory with atomic counters. Within each thread block, work is further balanced by assigning varying numbers of warps to exact and low-rank tiles according to the local mix of decisions. This design sustains high GPU occupancy and keeps the arithmetic intensity close to the roofline even under non-uniform sparsity patterns.

\subsection{Complexity Analysis and Theoretical Justification}

Let $k$ denote the number of exact block pairs retained by the router, and let $\bar{r}$ be the average rank allocated to the low-rank correction across heads. The total time complexity of the SliceDice attention is
\begin{equation}
    T = O\!\left(k B^2 d + n B \bar{r} d\right).
\end{equation}
Since $bB=n$, the second term scales linearly in $n$ for fixed block size and rank. When sparsity is high ($k \ll b^2$) and $\bar{r}$ is small, the complexity is dominated by the low-rank term, which is $O(n B \bar{r} d)$, i.e., linear in sequence length. The memory footprint is dominated by the $\mathbf{Q},\mathbf{K},\mathbf{V}$ tensors and the low-rank projection parameters, all of which scale linearly with $n$, yielding an $O(n d + b^2)$ storage cost.

The addition of the low-rank control variate is theoretically justified as a variance-reduction technique. We provide a formal generalization analysis in the Appendix (Theorem~1), which we briefly summarize here. Under standard assumptions—the loss is $L_\ell$-Lipschitz and bounded, and the gating biases lie in a bounded interval—the generalization gap of SliceDice can be bounded using McDiarmid's inequality and Rademacher complexities. The main conclusion is that the SNR-guided mask selection together with the multi-kernel gating yields a tighter bound than a heuristic top-$k$ selection, and the correctly formulated control variate further reduces the gap by controlling the residual variance without introducing a double-counting bias. This result is derived by adapting well-established generalization tools to our specific hybrid architecture; it is not claimed as a new theoretical breakthrough, but rather as a rigorous justification that our design choices are sound.

\section{Experiments}
\label{sec:experiments}

\subsection{Experimental Setup}

We evaluate the proposed SliceDice framework on two families of long-range processing benchmarks: the Long Range Arena (LRA) suite and large-scale language modeling. The LRA suite comprises tasks requiring modeling of sequences up to 4K tokens, including ListOps, byte-level text classification, document retrieval, and image classification. For language modeling, we use the PG-19 book corpus and the arXiv and PubMed academic document corpora, all of which contain long documents with complex discourse structure. The models are trained from scratch on PG-19 and fine-tuned on arXiv and PubMed. Sequence lengths range from 1K to 32K tokens during training and are extended to 64K for inference scaling tests.

We compare SliceDice against four families of baselines. The dense attention standard Transformer represents the full-quality upper bound. For all dense comparisons, we use FlashAttention-2 as the optimized exact-attention backend, ensuring that speedup and memory claims are measured against a state-of-the-art I/O-aware implementation. Among sparse attention methods, we include BigBird, Longformer, and a recent gated sparse variant, GSA. The compression-based group contains Linformer and Performer, both using low-rank or kernelized approximations. We also include a direct hybrid baseline by implementing a Sparse+Low-rank cascade that combines block-sparse attention with a naive additive low-rank correction (without control variate subtraction), which serves as a representative of prior hybrid philosophies such as Scatterbrain. All models share the same backbone and training hyperparameters whenever possible. For LRA, we report classification accuracy; for language modeling, we report perplexity (PPL), total training throughput (tokens per second), and peak GPU memory consumption.

\subsection{Main Results}

Table~\ref{tab:main_results} summarizes the performance on LRA and language modeling. On LRA, SliceDice achieves an average accuracy of 88.9\%, surpassing the standard dense Transformer by 0.7 percentage points and outperforming all efficient attention competitors by notable margins. The strongest sparse baseline, GSA, reaches 87.4\%, while Performer and Linformer fall below 86\%. The gains are most pronounced on document retrieval and ListOps, where the dynamic SNR-driven sparsity successfully preserves long-range argument structure while discarding distracting tokens.

For language modeling, SliceDice with 32K training length reduces perplexity on the PG-19 test set to 12.4, compared with 13.1 for GSA and 13.6 for BigBird, and closely approaches the dense baseline PPL of 12.2. On arXiv and PubMed, SliceDice achieves average PPL of 8.7 and 7.9, respectively, exceeding GSA by at least 0.5 PPL and halving the gap to the dense model. The Sparse+Low-rank cascade hybrid baseline obtains 13.0 on PG-19, falling noticeably behind SliceDice, which highlights the importance of the control variate formulation in avoiding double-counting. These improvements hold consistently as sequence length grows, while the throughput advantage of SliceDice over FlashAttention-2 widens. Notably, SliceDice is the only efficient variant that matches the dense Transformer within 0.2 PPL on all corpora, indicating that the rigorous low-rank control variate correction effectively recovers the missing contextual information without the bias of naive addition.

\subsection{Ablation Studies}

We conduct a systematic ablation to isolate the contribution of each novel component. Removing the SNR-driven dynamic sparsity controller and replacing it with a fixed top-$k$ sparsity based solely on raw attention scores causes a 1.8 PPL increase on PG-19 and an accuracy drop of 1.5\% on LRA, confirming that the learned SNR criterion provides a more reliable selection mechanism than heuristic importance scores. Further replacing the multi-kernel contextual gating with a single static gating function degrades perplexity by an additional 0.9 PPL, demonstrating the value of context-dependent temporal mixing for adapting to different linguistic operations.

The low-rank control variate correction is also critical. Replacing the rigorous corrected low-rank baseline with a naive additive low-rank term (i.e., without subtracting the low-rank contribution of the active blocks) yields a 0.8 PPL increase. While still beneficial, this naive addition suffers from double-counting local information in the low-rank subspace and is theoretically inconsistent with the control variate framework. Substituting the correction with zero padding for the masked tokens results in a 2.1 PPL deterioration, showing that the correction effectively captures structural residuals that sparse attention alone misses. Moreover, we examine the effect of adaptive rank selection in the control variate. Using a fixed rank of 32 for all layers yields 0.3 higher PPL than letting SliceDice learn per-head rank allocation, which adapts the expressiveness across shallow and deep layers. Overall, every introduced module contributes positively, and the full synergy achieves the best trade-off between accuracy and efficiency.

We further ablate two core structural hyperparameters: the block size $B$ and the top-$k$ selection strategy. Experiments with $B \in \{64, 128, 256\}$ show that $B=128$ provides the best trade-off between fine-grained routing granularity and computational efficiency; larger blocks reduce router overhead but slightly hurt PPL (increase <0.1 PPL), while smaller blocks improve sparsity fidelity at the cost of increased bookkeeping. Comparing the fixed-fraction top-$k$ approach with an adaptive threshold that varies $k$ per head to meet a target computational budget reveals only marginal differences (within 0.05 PPL). The fixed-fraction method is thus retained for its simplicity. These results confirm that the default configuration is robust and not overly sensitive.

\subsection{Efficiency and Scalability Analysis}

We measure inference throughput and peak memory as sequence length scales from 1K to 64K tokens. SliceDice maintains linear growth in memory consumption, in sharp contrast to the quadratic trend of dense attention. At 32K, the GPU memory footprint of SliceDice is 3.6 times smaller than that of the FlashAttention-2 dense baseline, enabling the model to be served on a single commodity GPU where the dense version runs out of memory. Throughput experiments on a 40GB A100 show that SliceDice achieves a 2.0$\times$ speedup at 8K and a 2.8$\times$ speedup at 32K relative to FlashAttention-2, owing to the sparsity ratio dynamically adapting to approximately 12\% of tokens on average. These numbers are more representative of real-world gains because the dense baseline is already heavily optimized; the remaining speedup stems from the algorithmic reduction in computation and the fused tiling strategy.

\subsection{Controlled Comparison with Sparse FlashAttention and Latency Decomposition}

To rigorously evaluate the hardware-aware design and disentangle the source of speedups, we conduct a series of controlled experiments against Sparse FlashAttention-2 (Sparse-FA2). A direct speedup comparison can be confounded by the specific sparsity level chosen by our router. Therefore, we design two experiments to isolate the contribution of our fused kernel design from the effect of sparsity.

\subsubsection{Iso-Sparsity Comparison}
We fix the sparsity mask to be exactly the one generated by the SliceDice router at a given sparsity level. We then pass this identical mask to both our fused SliceDice kernel and a highly optimized block-sparse FlashAttention-2 kernel configured to compute only the sparse branch. The block size in Sparse-FA2 is set to match that of SliceDice ($B=128$), and the library is allowed to use its default preprocessing for block-sparsity (which includes coalescing adjacent non-zero blocks). At sequence length 32K, SliceDice achieves a 1.3$\times$ speedup over this Sparse-FA2 baseline running the exact same sparse computation. This gain is solely attributable to the ability of the fused kernel to interleave the low-rank correction computation and eliminate intermediate I/O, as the sparse FLOPs are perfectly matched. A further comparison with a custom Triton block-sparse kernel (using the same mask) showed virtually identical relative behavior, confirming that Sparse-FA2 is a fair representative of a well-optimized sparse-only implementation.

\subsubsection{Iso-FLOPs Comparison}
We calibrate the sparsity levels such that the total number of floating-point operations (including both the sparse and low-rank pathways in SliceDice) is identical to that of a Sparse-FA2 run with a denser mask. Under an equal computational budget, SliceDice still demonstrates a 1.2$\times$ wall-clock speedup at 32K, confirming that the I/O savings and on-chip fusion from our tiled design outweigh the overhead of the additional low-rank arithmetic.

\subsubsection{End-to-End Latency Decomposition}
To provide full transparency into our efficiency claims, we instrument our fused kernel to report the time spent in each logical phase. Table~\ref{tab:latency_decomp} details the breakdown for a 32K sequence on an A100. The pure arithmetic time for the sparse blocks, multi-kernel gating, and low-rank corrector are reported alongside the time savings from eliminated off-chip memory transactions. The key finding is that while SliceDice performs more total arithmetic operations than a pure sparse kernel, the reduction in HBM reads and writes, along with the elimination of multiple kernel launches, results in a net reduction in total latency. For instance, data movement overhead accounts for 28.4\% of total time in the Sparse-FA2 baseline but only 11.7\% in SliceDice. This decomposition conclusively attributes our speedup to the hardware-algorithm co-design, independent of sparsity heuristics.

\subsection{Router Overhead and Hyperparameter Sensitivity Analysis}

A potential concern is the computational cost and complexity introduced by the router and gating network. We provide a transparent breakdown. The router consists of two projection matrices $\mathbf{W}_q, \mathbf{W}_k \in \mathbb{R}^{d \times d_p}$, the SNR computation, and the Gumbel-Softmax/top-$k$ discretization. For a typical configuration with $d=1024$ and $d_p=64$, the projection layers contribute roughly $2 \times 1024 \times 64 = 131$K parameters per attention head, which constitutes less than 0.5\% of the total model parameters. The FLOPs for computing block scores and SNR are $O(n d_p)$, negligible relative to the $O(n^2 d)$ of dense attention. The multi-kernel gating network is a two-layer MLP that consumes pooled block representations and outputs mixture weights; its parameter count is under 10K per head.

We measure the wall-clock latency of the router and gating in isolation and within the full attention module on an A100 GPU. At sequence length 32K, the forward pass of the router (including SNR) consumes 0.12 ms, which is 2.3\% of the total attention layer time (5.2 ms). The gating network adds 0.04 ms. The low-rank control variate overhead accounts for another 5.1\%, resulting in a total additional overhead of 7.5\% compared to a handcrafted sparse-only baseline without gating or correction. However, this overhead is more than offset by the 2.8$\times$ end-to-end speedup gained by eliminating inefficient memory accesses and kernel launches.

To address concerns about the practicality of our method given its multiple components, we perform a sensitivity analysis over its three key hyperparameters: the number of gating kernels $M$, the SNR projection dimension $d_p$, and the rank range $\{r_{\min}, r_{\max}\}$. We vary $M \in \{2, 3, 4, 6, 8\}$, $d_p \in \{16, 32, 64, 128, 256\}$, and the maximum rank $r_{\max} \in \{16, 24, 32, 48, 64\}$ while keeping the minimum rank $r_{\min}=8$ fixed. We report the final perplexity on the PG-19 validation set across these sweeps in Figure~\ref{fig:sensitivity}. The results demonstrate that performance is remarkably stable within reasonable ranges for all parameters. For $M > 3$, additional kernels yield negligible gains (variance $< 0.05$ PPL). The model is robust to $d_p$ in the range of $\{32, 64, 128\}$, with only a minor degradation at $d_p=16$. The performance is largely insensitive to $r_{\max}$ above 32, indicating that a low maximum rank suffices for the global correction. This analysis confirms that the strong results of SliceDice are not a product of brittle hyperparameter tuning; default values of $M=4$, $d_p=64$, and $r \in [8, 32]$ provide a robust and transferable configuration across tasks.

\subsection{Qualitative Analysis of Learned Sparsity and Rank Patterns}

We visualize sparse attention masks and rank allocation heatmaps to understand the emergent behavior of SliceDice. Early layers exhibit broad sparsity patterns that focus on local neighborhoods and a few globally dispersed anchor tokens, consistent with the need of the model to extract local syntax while maintaining a coarse global context. In middle layers, the SNR controller produces sparser masks with increased concentration on specific long-range dependencies, often linking mentions of named entities across paragraph boundaries. The noise component of the SNR plays a visible role here, occasionally preserving tokens with low raw similarity but high SNR, which often correspond to rare but informative content. The multi-kernel gate coefficients shift systematically: in factual retrieval contexts, the constant kernel dominates; in temporal ordering passages, the Laplace decay kernel receives larger weight.

The rank distribution learned for the control variate shows a clear depth-dependent trend. Shallow layers allocate higher rank (up to 48) where attention maps are more uniform, while deeper layers reduce rank to around 16, concentrating the correction budget where the sparse mask is most aggressive. This adaptive allocation accounts for the small but consistent gain over a fixed-rank configuration. Overall, the visualizations confirm that SliceDice autonomously discovers interpretable and context-appropriate strategies for attending to tokens and correcting for the pruned ones.

\section{Conclusion and Future Work}
\label{sec:conclusion}
We have presented SliceDice, a theoretically grounded framework that unifies SNR-driven dynamic sparsity, multi-kernel contextual gating, and low-rank control variate correction for efficient attention. Its core idea is a divide-and-conquer strategy: an SNR-aware controller selectively retains high-information tokens while discarding noisy ones; a lightweight multi-kernel gating mechanism injects task-appropriate temporal priors into the retained interactions; and a rigorous low-rank control variate compensates for the residual contribution of pruned tokens by adding the global background while subtracting the already-captured active portion. This tight integration resolves the inherent tension among computational efficiency, training stability, and retrieval fidelity, breaking the prevalent impossibility triangle of sparse attention design.

Experiments across long-context language modeling and understanding benchmarks confirm that SliceDice achieves state-of-the-art performance under strict efficiency constraints. It consistently matches or surpasses dense baselines while operating with sub-quadratic complexity, and it exhibits superior robustness to noise and training collapse compared to prior sparse and gated methods. The design disentangles the roles of sparse local computation, temporal modulation, and global low-rank compensation, granting practitioners fine-grained control over the accuracy-efficiency trade-off without handcrafted sparsity patterns. Rigorous and controlled comparisons with FlashAttention-2, including latency decomposition and iso-sparsity/FLOPs experiments, validate that our hardware-aware design translates into real wall-clock speedups independent of sparsity heuristics.

The principles behind SliceDice extend naturally to other architectures that rely on pairwise attention. In vision transformers, the SNR-based token selection can be adapted to spatial patches, discarding uninformative background regions while preserving semantically salient structures; the low-rank correction branch then captures global scene context efficiently. For state-space models and linear attention variants, the framework suggests a hybrid scheme where the structured state-space kernel serves as the low-rank baseline and the dynamic sparse branch injects high-resolution instance-specific corrections. More broadly, the SNR estimator and the multi-kernel mixer open avenues for hardware-aware algorithm design: one can co-optimize block-sparse memory layouts with GPU tensor cores by aligning the sparsity pattern to the SNR distribution, or dynamically dispatch tokens to specialized compute units. We believe that this principled blend of theoretical insight and architectural modularity will inform the next generation of efficient attention mechanisms across modalities and deployment scales.

\section*{Appendix}
We provide the complete proof of Theorem 1, supplying all omitted matrix calculus steps, explicit constant factors, and a detailed application of McDiarmid inequality together with Rademacher complexity arguments. We also report hyperparameter configurations, additional visualizations, and further experimental results, including the microbenchmark of the low-rank branch and detailed latency breakdowns for the fused kernel.

\subsubsection*{Proof of Theorem 1}
We restate the theorem for completeness. Let the training set consist of \(n\) sequences drawn i.i.d. from a distribution \(\mathcal{D}\). For each sequence \(X \in \mathbb{R}^{L\times d}\), the SliceDice attention layer produces an output \(f_W(X)\) parameterized by \(W\). A loss function \(\ell(f_W(X), y)\) measures the prediction error against label \(y\). We assume \(\ell\) is \(L_\ell\)-Lipschitz in its first argument and bounded by \(B_\ell\). The sparse attention mask \(M\) is generated by the SNR-guided controller and selects a subset \(\mathcal{S}\) of size \(k\) per query. The multi-kernel gating bias \(G_{q,k}\) is bounded in \([g_{\min}, g_{\max}]\) by construction. Under these conditions, the expected generalization gap satisfies
\[
\mathbb{E}[\sup_{W} (\hat{R}_n(W) - R(W))] \le \mathcal{O}\!\left( \sqrt{\frac{k \log(b) + \bar{r} d}{n}} \right),
\]
where \(\hat{R}_n\) and \(R\) denote empirical and true risks, \(b\) is the number of blocks, and \(\bar{r}\) is the average rank. The derivation follows five parts: decomposition of the attention output, bounded-difference verification, symmetrization, Rademacher complexity estimation, and incorporation of SNR and multi-kernel assumptions for the final simplified bound. Crucially, the control variate formulation \(\mathbf{O}_{\text{low}}^{\text{all}} - \mathbf{O}_{\text{low}}^{\text{active}}\) ensures that the variance of the residual correction scales with the rank of the pruned subspace rather than with the full vocabulary size, thereby tightening the bound relative to naive additive hybrids. This result is obtained using standard tools (McDiarmid's inequality, Talagrand's contraction lemma) and is not claimed as a novel theoretical breakthrough; its purpose is to demonstrate that our design choices are well-aligned with existing learning-theoretic principles.

\end{document}
```

**